% Generated draft: requires mathematical and source review.
\documentclass[11pt]{article}
\usepackage[margin=25mm]{geometry}
\usepackage{fontspec}
\setmainfont{Latin Modern Roman}
\usepackage{amsmath,amssymb,mathtools,unicode-math}
\setmathfont{Latin Modern Math}
\usepackage{enumitem,needspace,xurl,hyperref}
\setlength{\parindent}{0pt}
\setlength{\parskip}{5pt}
\setlength{\emergencystretch}{3em}
\newcommand{\N}{\mathbb{N}}
\newcommand{\Z}{\mathbb{Z}}
\newcommand{\Q}{\mathbb{Q}}
\newcommand{\R}{\mathbb{R}}
\newcommand{\C}{\mathbb{C}}
\newcommand{\F}{\mathbb{F}}
\newcommand{\A}{\mathbb{A}}
\newcommand{\PP}{\mathbb{P}}
\newcommand{\E}{\mathbb{E}}
\newcommand{\eps}{\varepsilon}
\newcommand{\dd}{\,\mathrm{d}}
\begin{document}
\section*{UnsolvedMath ID: 30004228}
\textbf{Source problem code:} catalogue entry
\subsection{Definitions and mathematical statement}
Let \(G=\Z/p\Z\), where \(p\) is an odd prime, and let \(n,k\) be integers with \(n>k\geq1\). Let \(S^n\) denote the \(n\)-sphere equipped with the free, orientation-preserving \(G\)-action in the source. Freeness means that \(gx=x\) implies \(g=0\) for every \(x\in S^n\); orientation preservation means that every group element preserves the orientation of the sphere. Let \(V\) be a real \(k\)-dimensional representation of \(G\). A map \(f:S^n\to V\) is equivariant if
\[
f(gx)=g f(x)\qquad(g\in G,\ x\in S^n).
\]
The source assumes a center map
\[
c:\operatorname{CO}(S^n)\longrightarrow S^n
\]
and the condition
\[
e(V)\notin\operatorname{Index}_{G}\bigl(SO(n+1);\F_p\bigr),
\]
where \(SO(n+1)\) is the group of real orthogonal \((n+1)\times(n+1)\) matrices of determinant one and \(\F_p\) is the field with \(p\) elements.

The supplied record does not define \(\operatorname{CO}(S^n)\), the axioms of a center map, the predicate ``\((k,\varepsilon)\)-pancake,'' or the convention for the equivariant index. It also does not specify the \(G\)-action on \(SO(n+1)\) needed for an equivariant index. The notation \(e(V)\) suggests an equivariant Euler class, but its definition and the ambient cohomology group for the displayed nonmembership condition are not supplied. These symbols and conditions are retained as source terminology; the formulation below depends on resolving them.

\emph{Requested statement.} For every choice of the data above satisfying the source's hypotheses, is it true that
\[
\forall\varepsilon>0\ \exists i_0\in\Z_{\geq0}\ \forall i\in\Z_{\geq0},\ i\geq i_0,
\]
there exists a convex equal-volume partition
\[
\mathcal P=\{P_1,\ldots,P_N\},\qquad N=p^{i+1},
\]
of \(S^n\), whose pieces belong to the domain of \(c\), such that
\[
\{gP_j:1\leq j\leq N\}=\mathcal P\qquad(g\in G),
\]
\[
f\bigl(c(P_j)\bigr)=0\in V\qquad(1\leq j\leq N),
\]
and every \(P_j\) is a \((k,\varepsilon)\)-pancake in the intended source sense? Here \(gP=\{gx:x\in P\}\), so invariance means that the action permutes the pieces. Equal volume means that all pieces have the same volume under the intended spherical volume measure. The threshold \(i_0\) may depend on \(\varepsilon\) and all the fixed data.

The source does not specify the geometric and boundary conventions for convex partitions, the metric used for pancakes, or the regularity requirements on the action, \(c\), and \(f\). In particular, no isometric-action hypothesis or continuity hypothesis has been added. Without the missing definitions and conventions, this is a quantified restatement of the requested assertion, not a fully resolved mathematical formulation.
\subsection{Short English statement}
For every positive \(\varepsilon\), the question asks whether all sufficiently large powers \(p^{i+1}\) occur as the number of pieces in a group-invariant convex equal-volume partition of the sphere, with every center mapped to zero by \(f\) and every piece a \((k,\varepsilon)\)-pancake. The center, pancake, and equivariant-index conventions require clarification. The supplied registry labels the problem open; this draft does not independently verify that status.
\subsection{Sources}
\begin{itemize}
\item Supplied local UnsolvedMath registry record, permanent ID \(30004228\), title ``Equivariant Pancake Partitions of Spheres,'' exact source problem code \texttt{catalogue entry}, in the collection ``Oberwolfach Reports Open Problems.'' The record gives proposed year 2019 and does not identify an individual proposer.
\item Supplied source citation: \emph{Geometric, Algebraic, and Topological Combinatorics} (2020), Oberwolfach Reports, pp.\ 2395--2472, DOI \texttt{10.4171/owr/2019/39}. Source URL and supplied literature-source link, labelled ``Oberwolfach report'': \url{https://doi.org/10.4171/owr/2019/39}. The supplied background describes Oberwolfach Reports as reports documenting mathematical workshops and their open-problem discussions. No precise internal location for this problem is supplied.
\item Literature assessment dated \texttt{2026-08-22}: the supplied curation labels the problem open and reports that the stronger \(p\)-adic Gromov--Borsuk--Ulam statement is presented as a theorem being sought, with no primary proof or disproof of that exact statement found in that review. This is supplied dated curation, not an independently verified literature assessment. Its stated checked source is the Oberwolfach report linked above.
\item The registry supplies the statement-status label \texttt{corrected\_verified}, explaining that notation, typography, and sentence structure were repaired while mathematical content was preserved. This is a supplied editorial qualification, not mathematical verification of this draft.
\end{itemize}
\end{document}
